% 兩類理由 × 兩個機關（1375、1377–1393）。三格灰階，決定勝負的一格用強調色。
\begin{tikzpicture}[qa,cell/.style={nbox,text width=52mm,minimum height=17mm}]
  \node[hdr] at (3.15,1.25) {結果面的理由（第肆部分）};
  \node[hdr] at (8.85,1.25) {程序面的理由（第伍部分）};
  \node[hdr,anchor=east] at (0.15,0) {立法機關};
  \node[hdr,anchor=east] at (0.15,-2.0) {法院};
  \node[cell] at (3.15,0) {代表結構使受影響者的處境得以體察（1377--1378）；議會辯論可直接處理道德問題（1384--1385）};
  \node[cell] at (8.85,0) {公平選舉與多數決，給予平等的發言權與平等的決定權威（1387--1389）};
  \node[cell] at (3.15,-2.0) {個案取向、文本取向、說明理由三項優點，經檢驗不成立或兩方均有（1379--1386）};
  \node[cell,kbox] at (8.85,-2.0) {法官的民主資格遠不及議員；大法官之間的簡單多數決欠缺公平性的論證（1391--1392）};
  \node[lbl,anchor=north] at (3.15,-3.05) {至多未有定論};
  \node[lbl,anchor=north] at (8.85,-3.05) {一面倒，不利於違憲審查};
\end{tikzpicture}
